\documentclass[11pt,a4paper]{article}
\usepackage[margin=25mm]{geometry}
\usepackage{amsmath,amssymb,amsthm}
\newtheorem{lemma}{Lemma}\newtheorem{proposition}[lemma]{Proposition}
\title{The orthogonality boundary in the Suzuki product calculation}
\author{Human Compression project: private F003 source audit}\date{}
\begin{document}\maketitle
This note concerns Spiga's fixed four-page preprint,
arXiv:2010.04080v1. It does not assert that its rank theorem is false.
We retain Wilson's characterization exactly as quoted on its page 2:
the symplectic matrices defining the Suzuki group preserve the displayed
semilinear product on orthogonal pairs. The parameter is
$q=2^{2e+1}$, $t=2^e$, so $2t^2=q$. The alternating form has
antidiagonal matrix $J$. All coordinates below are in that preprint's
basis, not an unannounced alternative matrix convention.

\begin{lemma}[The missing defect vector]\label{defect}
Let $D:V\times V\to V$ be alternating and additive in both variables,
with $D(\lambda u,\mu v)=\lambda^t\mu^tD(u,v)$. Then $D$ vanishes
whenever $f(u,v)=0$ if and only if there is a vector $r$ such that
\[
 D(u,v)=f(u,v)^t r\quad\hbox{for all }u,v.
\]
\end{lemma}
\begin{proof}
The four off-diagonal orthogonal basis pairs force
$D(e_1,e_2)=D(e_1,e_3)=D(e_2,e_4)=D(e_3,e_4)=0$.
The orthogonal pair $(e_1+e_2,e_3+e_4)$ then gives
$D(e_1,e_4)=D(e_2,e_3)=r$. Expanding by semilinearity yields
\[
D(u,v)=(u_1v_4+u_4v_1+u_2v_3+u_3v_2)^t r.
\]
This is the asserted formula. Its converse follows immediately.
\end{proof}

For a symplectic matrix $g$ define
$D_g(u,v)=gu\mathbin\bullet gv-g(u\mathbin\bullet v)$.
It has the hypotheses of Lemma~\ref{defect}, since the printed product
is alternating and semilinear. Thus Wilson's retained characterization
is equivalent to five vector identities: the four orthogonal basis
defects vanish and the two nonorthogonal basis defects are equal.
It does not require either of the latter defects to vanish.

In the symmetric-matrix notation $x=(a_{ij})$ used in the source, its
equations (13)--(20) set these two defects separately to zero. The
correct replacement is the following four equalities:
\begin{align*}
a_{12}a_{44}+a_{14}a_{24}&=a_{22}a_{34}+a_{23}a_{24},\\
a_{11}a_{24}+a_{12}a_{14}&=a_{12}a_{23}+a_{22}a_{13},\\
a_{13}a_{44}+a_{14}a_{34}&=a_{23}a_{34}+a_{33}a_{24},\\
a_{11}a_{34}+a_{13}a_{14}&=a_{12}a_{33}+a_{23}a_{13}.
\end{align*}
Here the inverse Frobenius automorphism removes the common exponent
$t$. In fact these are already off-diagonal entries of $xJx=J$ when
$x$ is symmetric. Equations (3)--(12) come from the orthogonal pairs
and are unaffected by this domain correction; (21)--(22) are valid
symplectic identities. The repeated $xe_1\bullet xe_3$ in the source's
line introducing $x(e_2\bullet e_4)$ is a separate index typo.

\begin{proposition}[A uniform intermediate counterexample]
In every field in the stated range the binary matrix
\[
x=\begin{pmatrix}
0&1&0&1\\1&1&1&0\\0&1&0&0\\1&0&0&0
\end{pmatrix}
\]
belongs to the source's Suzuki model and satisfies $xJx=J$, but is
neither $J$ nor diagonal. It contradicts the claimed classification
of $X=\{x\in\operatorname{Sz}(q):xJx=J\}$.
\end{proposition}
\begin{proof}
Multiplication gives $x^T=x$ and $xJx=J$. On the four orthogonal
basis pairs the defect vectors are zero; on $(e_1,e_4)$ and
$(e_2,e_3)$ both are $e_2$. These computations involve only zero and
one, so every Frobenius power has the same effect. Lemma~\ref{defect}
therefore proves the product identity on every orthogonal pair, and
the retained Wilson characterization gives $x\in\operatorname{Sz}(q)$.
The entries $a_{24}=0$, $a_{12}=a_{14}=1$ place it in source Case 3.2
and directly violate printed equation (14), whose two sides are one
and zero. This is a counterexample to the intermediate claim, not a
rank-four polytope marking.
\end{proof}

The source's concluding non-generation argument uses that intermediate
classification and is consequently unavailable. The accepted proper-face
argument and its subfield continuation recover part of the rank-four
exclusion by a different method. The case of two proper Suzuki subfield
face groups remains a separate principal obligation.
\end{document}
